\documentclass[instructions.tex]{subfiles}
\begin{document}
 
\chapter{De l'aumône.}
 
Y a-t-il une étroite obligation à faire l'aumône ? Comment faut-il la faire ? Quelle récompense doit-on en attendre ?
 
\primo Oui, on est obligé à faire l'aumône, pour quatre raisons. La première est ce que vous devez à Dieu. Il a droit de vous commander l'aumône, puisqu'il est le maître de vos biens; il vous la commande en effet. S'il vous fait riche, c'est en faveur des pauvres; il ne vous donne du bien que pour en user selon vos nécessités; et pour secourir les misérables. Le commandement de la charité du prochain, dit Jésus-Christ, est semblable au commandement d'aimer Dieu; il est donc aussi impossible d'être sauvé sans faire l'aumône, quand on peut la faire, qu'il est impossible d'être sauvé sans aimer Dieu. Celui qui a des biens de ce monde, dit saint Jean (remarquez ces paroles), il ne dit pas celui qui a beaucoup de biens, ou celui qui en a trop, mais celui qui a des biens et qui voit son frère dans la nécessité, s'il lui refuse du secours, comment peut-il dire qu'il aime Dieu ?\footnote{Qui habuerit substantiam hujus mundi, et viderit fratrem suum necessitatem habere, et clauserit viscera sua ab eo, quomodò caritas Dei manet in eo? 1. Joan. 3.}
 
Vous auriez beau former de grands désirs de plaire à Dieu, faire des pénitences, fréquenter les sacremens, dire à Dieu que vous l'aimez; c'est hypocrisie et mensonge, si vous êtes dur pour le pauvre.
 
Seconde raison. Ce que vous devez à Jésus-Christ. Il nous assure que tout ce que nous faisons au prochain, fût-il le dernier des hommes, il le tient fait à lui-même\footnote{Mihi fecistis. Matth. 25.}. C'est donc Jésus-Christ que vous assistez, ou que vous rebutez, lorsque vous assistez ou que vous rebutez le pauvre. C'est sur ce point que nous serons examinés au jugement. J'ai eu faim, dira-t-il aux réprouvés, et vous ne m'avez pas donné à manger; j'ai eu soif, et vous ne m'avez pas donné à boire; j'ai été nu, et vous ne m'avez pas vêtu; j'ai été malade, et vous m'avez abandonné; allez, maudits, au feu éternel\footnote{ Matth. 25.}.
 
Troisième raison. Ce que vous devez aux pauvres. Il est votre frère. Pour être misérable, est-il moins digne de charité ? Si vous avez plus de bien, vous ne l'avez pas mérité. Dieu pouvait vous faire naître dans l'indigence, et faire riche ce misérable qui attend votre secours. Quelque pauvre qu'il soit, il est l'enfant de Dieu, peut-être plus aimé de Dieu que vous. Il est cher à Jésus-Christ, racheté de son sang aussi bien que vous. Vous avez soin de vos animaux; vous avez pitié d'une bête qui souffre; estimez-vous moins un voisin pauvre, racheté du sang du Sauveur, qu'un vil animal ?
 
Ne méprisez donc jamais les gens pauvres, ils sont vos semblables, et Dieu les estime. Ecoutez, mes frères, dit saint Jacques, Dieu a choisi les pauvres; ils sont riches selon la foi, héritiers du royaume qu'il a préparé à ceux qui l'aiment. Si ils perdent le ciel, c'est leur faute. C'est aussi votre faute s'ils se damnent, parce que vous les abandonnez. Si les pauvres ont de grands vices, vous et les autres en avez encore de plus grands. Par le zèle que vous devez avoir pour leur salut, ne devez-vous pas empêcher leur damnation par vos libéralités ?
 
Quatrième raison. Ce que vous devez à vous-même. Si vous êtes riche, mais que vous pouvez devenir pauvre. Dieu ne manque pas de moyens pour vous réduire à la misère; peut-être dans peu de temps, en punition de votre dureté, y serez-vous réduit. Vous saurez alors ce que c'est que d'être abandonné. Vous êtes pécheur; vous avez besoin de la miséricorde de Dieu. Or vous devrez vous attendre à être jugé sans miséricorde, si vous manquez de charité pour vos frères\footnote{Judicium enim sine misericordiâ illi qui non fecit misericordiam. Jac. 2.}.
 
\secundo Comment faut-il faire l'aumône ? Avec joie, avec des sentimens de charité, en vue de Dieu, non par vanité. Faites-la abondamment, à proportion de vos biens. Tel donne quelques sous qui devrait donner des pièces d'or. Tel donne quelques morceaux de pain qui devrait donner des mesures et des sacs de grain. Donnez souvent, parce que tous les jours Dieu vous donne, et que les pauvres ont toujours besoin. Si vous avez peu, donnez peu; si vous n'avez rien, rendez d'autres services: l'aumône à tous ceux qui la demandent, dit Jésus-Christ\footnote{Omni autem petenti te, tribue. Luc. 6.}. Il vaut mieux la donner à vingt qui ne la méritent pas que de la refuser à un seul innocent.
 
Saint Grégoire donnait à tous. Ayant un jour logé dans sa maison Jésus-Christ, sous la figure d'un étranger, il dit ces paroles mémorables: Donnez à tous ceux qui vous demandent, de peur que celui à qui vous refuserez ne soit Jésus-Christ en personne.
 
Si vous ne pouvez donner à tous, préférez vos pauvres parens, les pauvres de votre lieu, vos pauvres débiteurs et vos sujets. Quant à certains mendians oisifs et robustes, la meilleure aumône est de leur procurer l'instruction et l'occupation. Elever de pauvres enfans vagabonds et abandonnés, est une charité des mieux placées. Faites l'aumône de votre bien; rendre ce qu'on a pris, c'est une restitution et non pas une aumône.
 
On doit multiplier ses aumônes; elles sont même plus indispensables dans les calamités. Un homme qui a des grains, des denrées, qui, pour profiter de la nécessité publique, les refuse ou les conserve pour les vendre cher, au pauvre peuple, est un homme détestable, dont l'avarice est en exécration à Dieu et aux hommes.
 
Je ne puis faire l'aumône, direz-vous, parce que je n'ai que le nécessaire, et point de superflu. Vous le dites, mais Dieu voit du superflu, et beaucoup de superflu dans vos dépenses. Avoir de grands biens, et n'avoir pas de superflu, est un grand crime. Un coeur bienfaisant a toujours de quoi donner, dit saint Léon; un coeur avare n'a jamais rien. On est pauvre dans le sein de l'abondance, quand on ne sait pas modérer sa cupidité ni borner ses dépenses. On est riche dans l'indigence, dit saint Martin de Brague, quand on sait se contenter de son état.\footnote{Qui sibi ipsi satis est, cum divitiis natus est. Prov. 28.} N'est-on pas assez riche quand on a le nécessaire pour passer cette vie et se disposer à l'autre ?
 
\tertio Quelle récompense nous procure l'aumône ? 1. Les biens temporels. Un verre d'eau, dit Jésus-Christ, donné en mon nom, ne sera pas sans récompense. L'avarice a ruiné un million de familles, l'aumône en a enrichi une infinité. Celui qui la fait, dit le Sage, ne sera jamais dans l'indigence\footnote{Qui dat pauperi, non indigebit; qui despicit pauperem sustinebit penuriam. Prov. 28.}. Si vous rejetez le pauvre, vos biens périront. C'est un grand trésor pour le temps de la nécessité que de faire l'aumône\footnote{(3) Tob. 4.}, disait Tobie à son fils.
 
2. Le pardon des péchés. Rachetez vos péchés par vos aumônes, disait un jour un prophète à un grand roi. L'aumone resiste au péché, dit le Sage, comme l'eau éteint le feu\footnote{Eccl.3}, pourvu toutefois qu'on quitte le péché, et qu'on change de vie. Vous ne pouvez faire de grandes pénitences. Donnez en aumône ce que vous avez de reste, dit Jésus-Christ, et tout sera purifié. Soyez assuré que Dieu vous fera miséricorde, si vous avez pitié du misérable.
 
3. L'aumône fait comme il faut, nous obtient une mort sainte et nous procure le ciel. Ah ! qu'on pourrait devant Dieu avec confiance, disait Tobie, lorsqu'on a fait l'aumône. Mon fils, l'aumône ne laissera pas tomber votre âme en enfer. Je n'ai jamais vu, disait saint Jérôme, qu'un homme charitable ait fait une fin malheureuse. Le pauvre fût-il un scélérat, l'aumône faite en vue de Dieu a toujours son mérite.
 
Concluons, avec saint Chrysostome, que sans le suffrage des pauvres vous ne serez jamais sauvé. S'ils plaidèrent notre cause au jugement, notre salut sera assuré; s'ils sont contre nous, notre damnation sera inévitable.
 
Un jour vous demanderez miséricorde à ceux à qui vous la refusez aujourd'hui; employez donc avec mérite, pendant qu'il est temps, des biens qui vous seront inutiles.
 
Il n'est point de pauvre qui ne puisse vous rendre riche dans l'éternité. Quand vous lui faites l'aumône, vous ne lui donnez que ce que vous devez à Jésus-Christ, et ce que vous devez à Dieu, ce que vous devez à votre frère, et ce que vous devez à vous-même. Quand vous la lui refusez dans sa nécessité, vous retenez ce qui lui est dû, vous devenez son meurtrier, dit saint Augustin\footnote{Si non pavisti, occidisti. S.Aug.}. Et vous faites tort à vous-même, puisqu'un coeur dur, selon l'oracle du Saint-Esprit, fera une mauvaise fin\footnote{cor durum habebit male in novissimo. Eccli. 3. 27.}.
 
\section{Résolutions}
 
\primo Je ne mépriserai jamais les pauvres, le fussent-ils par leur faute; 

\secundo je les assisterai selon mes facultés, avec joie, en vue de Dieu, non par ostentation et vanité. 


\prayer{O Jésus, inspirez-moi une tendre affection pour vos membres souffrans, et faites-moi la grâce de racheter mes péchés par mes aumônes.}
 
\end{document}
